\xiaosi

国内在自动驾驶规划领域发展迅速。清华大学MARS实验室提出DriveVLM\cite{tian2024drivevlm}，将视觉语言模型的链式推理能力引入驾驶规划。上海交通大学Thinklab团队构建了Bench2Drive\cite{bench2drive2024}闭环评测基准，包含44种交互场景和220条评测路线，推动了端到端方法的公平比较。

在产业应用方面，理想汽车已在量产车上部署视觉-语言-动作（VLA）架构；小鹏汽车的图灵芯片支持VLA与自主强化学习模型的运行。中国工信部于2024年发布《智能网联汽车准入管理》规定，要求L3及以上自动驾驶系统具备决策可追溯能力，这与欧盟AI法案的要求形成呼应，进一步推动了可解释自动驾驶的研究需求。

在可解释性研究方面，Atakishiyev等人\cite{atakishiyev2024xaiad}系统总结了可解释AI在自动驾驶领域的五类贡献；Zablocki等人\cite{zablocki2022xad}在IJCV综述中指出，黑箱模型缺乏因果理解能力是获取监管信任的核心障碍。国内外学者普遍认为，\textbf{内在可解释架构}（如NCP的结构化稀疏设计）比事后解释方法（如Grad-CAM）更适合高风险的自动驾驶场景。

综合来看，当前自动驾驶规划领域存在"性能-效率-可解释性"三角困境：高性能模型参数量大、不可解释；轻量模型性能不足；可解释方法多为事后分析，无法提供结构性保证。NCP以其极少参数、内在可解释性和生物启发的稀疏结构，为打破这一困境提供了可能。
